\chapter{Geometric Extraction of the Local Navigation Reference}
\label{cap:geometric_reference}

This chapter describes the two geometric extractors implemented in
\texttt{pf\_\allowbreak straight\_\allowbreak prior3\_\allowbreak single.py}
and
\texttt{pf\_\allowbreak straight\_\allowbreak prior4\_\allowbreak single.py}.
They are referred to as Method~3 and Method~4, respectively. Both convert
a segmented corridor mask into a local navigation reference by extracting
intervals in horizontal strips, associating their boundaries from bottom
to top, and fitting the left and right edges separately. The reference
is the mean of the two fitted edge lines, rather than a line fitted
directly to interval midpoints.

The methods share mask preparation, interval extraction, and residual-based
line reweighting. They differ in the support used to select the initial
interval, the association cost, the handling of missing observations,
and the checks applied before accepting the output. This chapter explains
the active implementations and their default parameters; comments that
describe inactive or intended behavior are not treated as implemented
operations. Each image is processed independently. Temporal filtering,
conversion to velocity commands, and closed-loop evaluation belong to
the robotic integration stage.

\section{Input representation and edge-line models}
\label{sec:c4_input}

Both functions expect a two-dimensional label mask $M(y,x)$ of height $H$
and width $W$. Image coordinates increase to the right in $x$ and downward
in $y$. The argument \texttt{path\_class} identifies the class from which
corridor intervals are extracted; its default is 1 in both files. A
binary input can therefore use 1 for the target region and 0 elsewhere.
For a multiclass segmentation mask, this argument must match the actual
label encoding, or the mask must first be converted appropriately. The
extractors do not infer the class mapping or use a separate plant class.
A target-region label is not, by itself, a guarantee of traversable space.

The left and right boundaries are represented by
\begin{equation}
    x_L(y)=m_Ly+q_L, \qquad x_R(y)=m_Ry+q_R.
    \label{eq:c4_edges}
\end{equation}
Treating $x$ as a function of $y$ permits vertical image lines without an
infinite slope. The center reference is
\begin{equation}
    x_C(y)=\frac{x_L(y)+x_R(y)}{2}=m_Cy+q_C,
    \qquad m_C=\frac{m_L+m_R}{2},\quad q_C=\frac{q_L+q_R}{2}.
    \label{eq:c4_line}
\end{equation}
Thus, both methods assume locally straight image-space boundaries.
The models do not reconstruct terrain or estimate a metric trajectory.
Their residuals and width estimates are in pixels, not lateral distances
in meters. The two edges need not be parallel: their separation can
change linearly with image height.

At each image height, the reference lies halfway between the two
predicted boundaries. This construction is not generally equivalent to
fitting the observed interval midpoints: the two edge fits can exclude
different border-touching samples and assign different residual weights.
Keeping the boundaries separate allows each visible side to contribute
its own observations, although both sides must have enough samples.

The shared processing stages are:
\begin{enumerate}
    \item select a lower region of interest and clean the target-class mask;
    \item extract and merge projected intervals in horizontal strips;
    \item select a seed using the method-specific support and priority rules;
    \item associate interval boundaries upward through the image;
    \item fit the non-clipped left and right boundaries independently;
    \item calculate quality indicators and rasterize the center reference,
    with method-specific validation and output rules.
\end{enumerate}

\section{Region of interest and mask preparation}
\label{sec:c4_roi}

The ROI begins at $y_0=\lfloor\rho H\rfloor$ and extends to the bottom
of the image. The default starting ratio is $\rho=0.20$ for Method~3
and $\rho=0.30$ for Method~4. The binary target mask is
\begin{equation}
    B(y,x)=\mathbf{1}\{M(y_0+y,x)=\texttt{path\_class}\},
    \qquad 0\leq y<H_R=H-y_0.
\end{equation}
Both implementations apply opening with a $5\times5$ elliptical kernel,
then closing with a $9\times9$ rectangular kernel. Connected components
with fewer than 500 pixels are removed; all sufficiently large components
are retained, not only the largest one.

Let $\widetilde B$ denote the cleaned binary ROI mask. It supplies both
interval extraction and local support in the new methods. Unlike the
previous midpoint-based implementation, these support calculations do
not use the original multiclass mask to compute separate soil and plant
ratios. Method~3 additionally consults the original input labels when
assessing the rasterized centerline. Method~4 does not perform that
centerline-label assessment.

Morphology and interval merging can bridge interruptions or change region
connectivity. Consequently, support in $\widetilde B$ measures agreement
with the processed mask, not independently verified scene geometry.
Kernel sizes, component areas, and distance thresholds are in input-mask
pixels and depend on image resolution and camera configuration.

\section{Horizontal strips and candidate intervals}
\label{sec:c4_candidates}

\subsection{Strip construction and column projection}
\label{subsec:c4_strips}

Both methods use $N=24$ strips by default, indexed from bottom to top.
Their height is $h=\max(1,\lfloor H_R/N\rfloor)$. For strip index $s$,
the ROI-relative bounds are
\begin{equation}
    y_s^- = \max(0,H_R-(s+1)h),\qquad
    y_s^+ = \min(H_R,H_R-sh),
\end{equation}
with the lower image row bound $y_s^+$ excluded. Empty strips are skipped.
The representative full-image coordinate is
$y_s=y_0+(y_s^-+y_s^+)/2$, following the code's bound-based convention.
Integer division may leave a remainder above the constructed strips.

For example, with $H_R=157$ and $N=24$, the strips cover the lowest
$24\times6=144$ ROI rows. The upper 13 ROI rows do not supply candidates;
the ROI and the region actually sampled by the strips are therefore not
identical.

Within each strip, the cleaned mask is summed vertically. A column is
active whenever its sum exceeds zero, so even one target pixel activates
that column. Contiguous active columns form raw intervals. Neighboring
intervals are merged when the number of inactive columns between them
is at most \texttt{max\_interval\_gap}, whose default is 30 pixels.
For a merged interval $[\ell_i,r_i]$, the width is
\begin{equation}
    d_i=r_i-\ell_i+1.
    \label{eq:c4_interval}
\end{equation}
Intervals narrower than 8 pixels are discarded. There is no fixed maximum
interval width in either new method. Instead, an interval touching both
image margins is rejected if the mean of the entire cleaned binary strip
is below 0.70. With \texttt{border\_margin}=2, the flags are
\begin{equation}
    t_{L,i}=\mathbf{1}\{\ell_i\leq2\},\qquad
    t_{R,i}=\mathbf{1}\{r_i\geq W-1-2\}.
\end{equation}
These flags indicate possible truncation by the image boundary, rather
than verified physical corridor edges.

Column projection merges evidence across different vertical locations;
a merged interval can also contain horizontal gaps. Neither its width
nor its endpoints prove that its rectangular interior is entirely target
class. This matters when interpreting intervals around occlusions or
fragmented segmentation.

\subsection{Method-specific local support}
\label{subsec:c4_support}

Method~3 computes central support $u_i$ in a band around the rounded
interval midpoint. With default half-width 5, the band contains up to
11 columns and spans the strip height. Its value is the mean of
$\widetilde B$ inside the band, clipped at the image boundaries.
Candidates must satisfy $u_i\geq0.27$ for both seed selection and
subsequent association.

This central band checks whether the projected interval contains target
pixels near its midpoint, which is useful when merging has joined two
regions across a gap. It does not check the whole interval or establish
that both endpoints are reliable.

Method~4 instead computes $u_{L,i}$ and $u_{R,i}$ in separate bands
centered on $\ell_i$ and $r_i$. The effective half-width is
$\max(2,\texttt{local\_band\_half\_width})$, equal to 5 by default.
These means quantify target-pixel occupancy around each boundary; they
are not boundary-detection probabilities. Because a band straddles a
boundary, it need not be entirely target class even for a clean interval.

In Method~4, edge support ranks seed candidates and contributes to the
final quality score. It is not an acceptance threshold and does not
weight the line fit. In particular, \texttt{min\_edge\_support}=0.20
is present in the signature but unused in the active implementation.

\section{Initialization and bottom-up association}
\label{sec:c4_selection}

\subsection{Method 3: lowest supported seed}
\label{subsec:c4_seed3}

Method~3 searches the first six constructed strips from the bottom, or
all available strips if fewer exist. It selects the lowest strip with at
least one candidate satisfying $u_i\geq0.27$. Within that strip, it
minimizes distance between the interval midpoint and $(W-1)/2$; ties
favor the widest interval. There is no consecutive-run eligibility test.

Let $\ell_{\mathrm{prev}},r_{\mathrm{prev}},d_{\mathrm{prev}}$ denote
the last accepted interval. In each strip above the seed, define
\begin{equation}
    \delta_i=\max\bigl(|\ell_i-\ell_{\mathrm{prev}}|,
                            |r_i-r_{\mathrm{prev}}|\bigr),\qquad
    D_3=\max(42,0.5d_{\mathrm{prev}}).
\end{equation}
Among candidates with adequate central support and $\delta_i\leq D_3$,
the method minimizes
\begin{equation}
    J_{3,i}=\delta_i+0.1|d_i-d_{\mathrm{prev}}|.
    \label{eq:c4_cost3}
\end{equation}
A strip without an admissible candidate is skipped, leaving the previous
interval unchanged. There is no maximum number of consecutive missed
strips, and the jump threshold is not increased to account for a longer
vertical gap. Strips below the selected seed are not revisited.

\subsection{Method 4: edge-support seed and gap-limited association}
\label{subsec:c4_seed4}

Method~4 pools all candidates in the first six constructed strips. It
maximizes $\min(u_{L,i},u_{R,i})$, then interval width. Unlike Method~3,
it does not prioritize proximity to the image center or explicitly choose
the lowest eligible strip. Exact ties retain the first candidate in the
bottom-up, left-to-right traversal. No minimum-support threshold is
applied, so even a weakly supported candidate can initialize tracking.

The boundary displacement $\delta_i$ is defined as above, but the jump
limit and association cost are
\begin{equation}
    D_4=\max(28,0.5d_{\mathrm{prev}}),\qquad
    J_{4,i}=\delta_i+0.4|d_i-d_{\mathrm{prev}}|.
    \label{eq:c4_cost4}
\end{equation}
Only the displacement limit determines admissibility; edge support does
not enter this cost or an additional tracking gate. Compared with
Method~3, the default lower bound on allowable displacement is smaller
and width change has a larger cost coefficient. For wide intervals,
the width-dependent term can dominate the jump limit in either method.

A counter records consecutive strips without a compatible candidate and
resets after each successful association. With
\texttt{max\_gap\_strips}=3, tracking stops on the fourth consecutive
miss. Stopping does not directly return \texttt{lost}: the accumulated
points proceed to fitting if at least three are available. The maximum
miss count, including the miss that triggers stopping, enters the quality
score. The previous interval remains unchanged during gaps.

The stopping rule and gap penalty have different effects: three misses
are tolerated for continued association, but already reduce the gap
score to zero. A fourth miss stops the search without automatically
rejecting a reference supported by the observations collected earlier.

Both methods perform greedy spatial association within one frame, not
temporal tracking. Each accepted interval contributes two boundary
observations, and at least three accepted intervals are required before
the edge-fitting stage.

\section{Separate robust fitting of the corridor edges}
\label{sec:c4_fitting}

\subsection{Excluding image-boundary observations}
\label{subsec:c4_edge_samples}

The left fit uses only accepted intervals not flagged as touching the
left image margin. The right fit independently excludes intervals
flagged as touching the right margin. At least two observations per edge
must remain. The two fits can therefore use different strips and span
different vertical ranges. A clipped boundary is not replaced by a
one-sided centerline estimate: if either edge lacks sufficient samples,
both implementations return an empty output with \texttt{lost} status.

The retained samples are interval endpoints from the cleaned column
projection. They are not points obtained by a dedicated edge detector.
Excluding border-touching samples reduces the influence of image
truncation, but does not establish the correctness of the remaining
boundaries.

\subsection{Residual-based reweighting}
\label{subsec:c4_robust_fit}

Both methods start each edge fit with uniform weights. There is no
semantic, interval-width, or vertical-position base weighting. For an
edge sample $(x_i,y_i)$, four reweighting iterations are performed by
default, with scale $\sigma=18$ pixels:
\begin{equation}
    e_i^{(t)}=x_i-m^{(t)}y_i-q^{(t)},\qquad
    w_i^{(t+1)}=\frac{1}{1+(e_i^{(t)}/\sigma)^2}.
    \label{eq:c4_reweight}
\end{equation}
Each iteration calls \texttt{numpy.polyfit} with degree one and updated
weights. Its weights multiply unsquared residuals
\cite{numpyPolyfitDocumentation}; the fitted objective at that iteration is
\begin{equation}
    (m,q)=\mathop{\arg\min}_{m,q}
        \sum_i \left[w_i(x_i-my_i-q)\right]^2.
    \label{eq:c4_objective}
\end{equation}
Thus, the effective coefficient of each squared residual is $w_i^2$,
not $w_i$. The implementation should not be interpreted as a different
standard robust-loss estimator solely from the form of its weight rule.
Weights are recomputed from the current model rather than accumulated
across iterations; no convergence-based stopping criterion is used.
Method~3 additionally checks for negligible total weights in its helper.

The final residual for each retained sample is its absolute horizontal
deviation from the corresponding fitted edge. Define $\bar e$ as the
mean over the concatenated left and right residual arrays, and
\begin{equation}
    F=\operatorname{clip}(1-\bar e/\sigma,0,1).
    \label{eq:c4_fit_score}
\end{equation}
Edges with more retained samples contribute more terms to this mean.
Reweighting reduces the influence of inconsistent points but cannot
correct an erroneous association chain that supports a coherent wrong
line. It does not impose parallel edges or estimate genuine path curvature.

\section{Quality indicators, validation, and image-space output}
\label{sec:c4_output}

\subsection{Method 3: centerline support after rasterization}
\label{subsec:c4_output3}

Method~3 constructs the mean edge line and rasterizes it from the top
of the highest accepted strip to the exclusive lower bound of the seed
strip. It does not extend the line below an elevated seed. Within this
range, it fills rows across skipped strips as well. Predicted horizontal
positions are clipped to $[0,W-1]$ and converted to integers; dilation
with a $3\times3$ elliptical kernel thickens the line. Dilation may extend
the nonzero output slightly beyond the pre-dilation row range.

Before dilation, the original input labels are checked in a band of
half-width 3 around each rasterized point. Let $P$ be the fraction of
pixels equal to \texttt{path\_class} across these clipped bands, and
$T$ the number of accepted intervals divided by the number of constructed
strips. The score is
\begin{equation}
    C_3=0.50T+0.30P+0.20F.
    \label{eq:c4_confidence3}
\end{equation}
The status is \texttt{lost} if $C_3<0.38$ or $T<0.25$,
\texttt{degraded} if neither condition holds but $C_3<0.60$, and
\texttt{ok} otherwise. These checks follow rasterization: a result marked
\texttt{lost} can still contain a nonempty centerline matrix. No explicit
edge-order or positive-width test is applied.

The diagnostic dictionary includes both edge models, the centerline
model, row bounds, residual and support indicators, clipped-point counts,
and the cleaned full-image ROI mask. Early selection or fitting failures
instead return a zero matrix and a reason string, with additional details
in some branches.

\subsection{Method 4: geometric checks and a composite score}
\label{subsec:c4_output4}

Method~4 validates the fitted boundaries before constructing the output.
At 16 evenly spaced heights covering the combined range of retained
edge samples, it evaluates
\begin{equation}
    D(y)=x_R(y)-x_L(y).
\end{equation}
All sampled widths must be strictly positive. The median of these
sampled widths is returned as \texttt{width}. This check concerns the
sampled observation range, not any extrapolation below it.

A separate indicator uses all accepted intervals, including those with
border-touching endpoints. For $K$ intervals in association order, let
$z_j=r_j-\ell_j$; unlike extraction width $d_j$, this quantity omits the
inclusive-endpoint $+1$. The width-irregularity indicator is
\begin{equation}
    E_D=\frac{Q_{0.90}\left(\{|z_{j+2}-2z_{j+1}+z_j|\}_{j=1}^{K-2}\right)}
                   {\max(\operatorname{median}_j|z_j|,1)},
    \qquad A=\operatorname{clip}(1-E_D/0.30,0,1),
    \label{eq:c4_width_score}
\end{equation}
where $Q_{0.90}$ is the 90th percentile. These are second differences in
sample order, not derivatives normalized by vertical distance. If strips
are skipped, sample spacing is unequal; even a width varying linearly
with image height need not yield zero second differences. The indicator
therefore measures an implementation-specific irregularity, not a
perspective-invariant physical curvature.

Let $S$ be the mean of the left support values used by the left fit and
the right support values used by the right fit, pooled together. With
maximum consecutive miss count $g_{\max}$ and allowed gap $g_a=3$,
\begin{equation}
    G=\operatorname{clip}(1-g_{\max}/\max(g_a,1),0,1),\qquad
    C_4=0.25T+0.20S+0.25F+0.20A+0.10G.
    \label{eq:c4_confidence4}
\end{equation}
The result is \texttt{lost} if the positive-width check fails,
$\bar e>2\sigma$, or $C_4<0.38$. Otherwise, it is \texttt{degraded}
for $C_4<0.60$ and \texttt{ok} for $C_4\geq0.60$.
There is no separate $T<0.25$ rejection rule in Method~4.

A result marked \texttt{lost} returns a zero output matrix before
rasterization, with no centerline model in the post-fit rejection branch.
For an accepted result, the mean edge line is rasterized from the top of
the highest accepted strip to the bottom of the image, then dilated as
in Method~3. This can extrapolate below the seed. It does not check target
labels along the rendered centerline, and the edge-order check does not
cover the entire rendered range when that range exceeds the sampled
observation range.

Diagnostics report the edge and centerline models, mean residual,
support, tracking fraction, median fitted width, width irregularity,
maximum gap, row bounds, and cleaned ROI mask. Earlier failures return
\texttt{lost} with reason strings such as \texttt{no\_seed\_candidate},
\texttt{tracking\_failed}, or \texttt{insufficient\_edge\_points}.

\subsection{Interpretation of status and confidence}
\label{subsec:c4_failures}

The two confidence values are heuristic scores with different components;
they are not calibrated probabilities and should not be compared as if
an identical numerical value represented identical reference quality.
Method~3 assesses original-mask support along the centerline, whereas
Method~4 combines processed-mask edge support with geometric and gap
indicators. Neither verifies obstacle-free passage.

A downstream controller must inspect the returned status. In particular,
a nonempty Method~3 output does not necessarily denote accepted guidance.
Both methods can also clip an out-of-image prediction onto an image
border. Neither function specifies a robot stopping or recovery policy.
Malformed inputs, unsuitable parameter values, or numerical exceptions
are not comprehensively handled by the explicit failure branches.

\section{Default configurations and comparison}
\label{sec:c4_discussion}

Table~\ref{tab:c4_defaults} distinguishes shared defaults from
method-specific settings. These are the defaults of the supplied files,
not a statement that a controller uses them or that they are
experimentally optimal. Overrides and mask encodings must be recorded
for each experiment.

\begin{table}[htbp]
    \centering
    \small
    \caption{Default configurations of the two edge-based extractors.
    Distances and areas are expressed in input-mask pixels.}
    \label{tab:c4_defaults}
    \begin{tabularx}{\textwidth}{@{}X l l@{}}
        \toprule
        Setting & Method 3 & Method 4 \\
        \midrule
        Target-class identifier & 1 & 1 \\
        ROI starting ratio & 0.20 & 0.30 \\
        Number of strips & 24 & 24 \\
        Opening / closing kernels & $5\times5$ / $9\times9$ & Same \\
        Minimum component area & 500 & 500 \\
        Minimum interval width & 8 & 8 \\
        Maximum merged gap & 30 & 30 \\
        Image-border margin & 2 & 2 \\
        Full-strip occupancy gate & 0.70 & 0.70 \\
        Local support-band half-width & 5 & 5 \\
        Seed search window & 6 strips & 6 strips \\
        Central support threshold & 0.27 & Not used \\
        Minimum edge support threshold & Not used & Inactive parameter \\
        Jump-limit lower bound & 42 & 28 \\
        Width-change cost coefficient & 0.1 & 0.4 \\
        Consecutive misses allowed & No fixed limit & 3 \\
        Residual scale / iterations & 18 / 4 & 18 / 4 \\
        Centerline assessment half-width & 3 & Not used \\
        Width-error score scale & Not used & 0.30 \\
        Lost / acceptable score thresholds & 0.38 / 0.60 & 0.38 / 0.60 \\
        Minimum tracking fraction gate & 0.25 & Not used \\
        Output dilation kernel & $3\times3$ & $3\times3$ \\
        \bottomrule
    \end{tabularx}
\end{table}

The principal difference is not the line model but the selection and
validation surrounding it. Method~3 favors a supported interval near the
image center in the lowest eligible strip, accepts arbitrarily long
association gaps, and evaluates the rendered line against the original
mask. Method~4 favors joint support of the two boundaries across the
bottom search window, limits consecutive gaps, penalizes width changes
more strongly, and rejects outputs failing its geometric or score checks.
These choices can affect different failure cases, but they do not establish
that one method is more accurate or faster without evaluation.

Both methods can merge separated regions, follow an incorrect but
consistent interval chain, or infer a line through an unobserved gap.
A wide or well-supported interval is not necessarily the intended row.
The two-line representation retains boundary information separately,
but requires sufficient non-clipped samples on both sides. Independent
edge fits may be supported over different height ranges, and their mean
can include extrapolation of one or both boundaries.

Evaluation should separately assess interval association, edge fitting,
centerline accuracy, rejection behavior, and closed-loop motion.
A controlled comparison should use the same input masks and first align
shared settings such as ROI and strip count, so differences are not
attributed solely to algorithmic rules when input regions also differ.
Useful component comparisons include seed priority, width-change cost,
gap stopping, and post-fit validation. The descriptions here specify
implemented behavior, not results of those comparisons.

\section{Effect of strip height and number}
\label{sec:c4_strip_height}

The strip count controls vertical sampling and the maximum number of
interval observations. More strips provide thinner vertical regions and
potentially more points per boundary fit. Taller strips aggregate more
mask pixels but mix evidence from different image heights, and column
projection can conceal interruptions. Increased strip height does not
guarantee more reliable boundary estimates or a numerically better fit.

For a $224\times224$ input, Method~3 has $y_0=44$ and $H_R=180$;
Method~4 has $y_0=67$ and $H_R=157$. Table~\ref{tab:c4_strip_tradeoff}
uses the exact integer-division rule for these illustrative dimensions,
rather than assuming the same ROI height for both methods. At $N=24$,
the strip heights are 7 and 6 pixels, respectively. A fit can use fewer
than $N$ points because of seed location, missed strips, early stopping,
or border-point exclusion; the two edges can have different sample counts.

\begin{table}[htbp]
    \centering
    \small
    \caption{Strip heights for an illustrative $224\times224$ mask and
    the default ROI ratios. The point count is an upper bound per edge,
    not a guaranteed number of fitted samples.}
    \label{tab:c4_strip_tradeoff}
    \begin{tabular}{@{}c c c c@{}}
        \toprule
        Strips $N$ & Method 3 height & Method 4 height & Max. edge samples \\
        \midrule
        24 & 7 px & 6 px & 24 \\
        20 & 9 px & 7 px & 20 \\
        16 & 11 px & 9 px & 16 \\
        12 & 15 px & 13 px & 12 \\
        8 & 22 px & 19 px & 8 \\
        \bottomrule
    \end{tabular}
\end{table}

Changing $N$ also changes the vertical extent of a fixed six-strip seed
window and, in Method~4, the extent of a three-strip tolerated gap. These
settings should therefore be considered together. The width-irregularity
indicator depends on the sampling sequence, while tracking fractions
use the number of constructed strips as their denominator. Selecting
$N$ requires evaluating boundary and centerline errors, accepted-point
counts, rejection rates, and processing time under the intended image
resolution and field conditions. No optimal strip count is inferred from
the implementation alone.
